\documentclass[conference]{IEEEtran}
\IEEEoverridecommandlockouts
\usepackage{cite}
\usepackage{amsmath,amssymb,amsfonts}
\usepackage{algorithm}
\usepackage{algorithmic}
\usepackage{graphicx}
\usepackage{textcomp}
\usepackage{xcolor}
\usepackage{booktabs}
\usepackage{url}
\usepackage{microtype}
\usepackage{balance}
\renewcommand{\algorithmicrequire}{\textbf{Input:}}
\renewcommand{\algorithmicensure}{\textbf{Output:}}
\def\BibTeX{{\rm B\kern-.05em{\sc i\kern-.025em b}\kern-.08em
    T\kern-.1667em\lower.7ex\hbox{E}\kern-.125emX}}
\begin{document}
% =============================================================================
\title{Computational Entropy: A Framework for Quantifying the Inference
Energy Cost of Prompt Semantic Instability in Large Language Models}
\author{
  \IEEEauthorblockN{Febin K Renu}
  \IEEEauthorblockA{
    \textit{Division of Computer Science and Engineering}\\
    \textit{Karunya Institute of Technology and Sciences}\\
    Coimbatore, Tamil Nadu, India\\
    febink@karunya.edu.in}
  \and
  \IEEEauthorblockN{G Naveen Sundar}
  \IEEEauthorblockA{
    \textit{Division of Computer Science and Engineering}\\
    \textit{Karunya Institute of Technology and Sciences}\\
    Coimbatore, Tamil Nadu, India\\
    naveensundar@karunya.edu}
}
\maketitle
% =============================================================================
\begin{abstract}
Prompt engineering research optimises overwhelmingly for output
quality; almost none of it asks what a prompt costs to process. This
paper contributes a \emph{framework} for asking that question
precisely --- not a proof that the cost exists. We propose
\emph{Computational Entropy} ($S_c$), a construct drawn from
semantic-entropy clustering, as the theoretical basis for a practical,
single-pass \emph{Semantic Instability Index} (SII); a nine-class
\emph{Mutation Engine} for generating semantically graded prompt
variants; and the \emph{Prompt Entropy Coefficient} (PEC), a Spearman
correlation between SII and measured energy-per-token (EPT). We then
apply this framework at three levels of measurement fidelity, and
report what each level actually found. Under a CPU timing-proxy
protocol with a synthetic energy term, PEC $=0.408$ ($p<0.001$,
$N=540$) and mutation condition explains $\eta^2=0.157$ of EPT
variance. Removing the synthetic term and measuring a real, from-scratch,
untrained Transformer on real CPU timing reproduces the direction
(PEC $=0.252$) but shows it is substantially mediated by mutation-induced
prompt length rather than semantic content. The decisive test replaces
both the untrained model and the timing proxy: two real, open-weight,
instruction-tuned models (Phi-3.5-mini, Gemma-2-2B), served locally
and measured with actual NVIDIA NVML GPU-power integration, with an
explicit length-matched control for every mutated prompt. On this real
hardware, the SII--EPT correlation is statistically indistinguishable
from zero for both models, before and after length matching
($|\rho|\le0.07$, all $p>0.4$), and a 200-draw coefficient-sensitivity
sweep confirms this null is stable rather than fragile. We report this
outcome as the paper's main empirical finding, not a footnote: the
SII/PEC/Mutation-Engine framework and its three-tier validation
protocol are reusable contributions in their own right, but the
evidence gathered here does not support treating prompt semantic
instability as a demonstrated energy cost in real deployed models. We
discuss why proxy and untrained-model measurements can manufacture an
association that real hardware does not confirm, what that implies for
energy claims elsewhere in the Green AI literature, and which two
narrower, model-specific findings keep the underlying question open
rather than closed. Code, corpus, mutation logic, and full measurement
logs for all three stages are made available for independent
verification.
\end{abstract}
\begin{IEEEkeywords}
LLM inference energy; prompt engineering; semantic instability;
computational entropy; energy-per-token; Green AI; replication;
NVML power measurement
\end{IEEEkeywords}
% =============================================================================
\section{Introduction}
Inference energy is usually treated as a function of model size,
hardware, and batch configuration~\cite{strubell2019energy,
patterson2022carbon,samsi2023words,desislavov2023trends}, and for
widely deployed models inference now exceeds training's share of
lifetime energy~\cite{wu2022sustainable,faiz2024llmcarbon}. Almost all
resulting efficiency work --- quantisation, speculative decoding,
prompt caching, carbon-aware scheduling~\cite{chien2023reducing,
dodge2022measuring} --- operates on the serving infrastructure and
takes the prompt itself as a fixed, unexamined input. Meanwhile prompt
engineering has grown into a discipline of its own
\cite{sahoo2024systematic}, almost entirely oriented toward output
quality. We ask a narrower, testable question that sits at the
intersection of the two literatures: is a prompt's semantic
instability --- incoherence, ambiguity, internal contradiction ---
statistically associated with the energy cost of generating a response
to it? A prompt containing a logical contradiction forces a model to
reconcile irreconcilable constraints while still producing an answer;
whether that extra burden shows up as measurable energy is an
empirical question, and this paper is an attempt to answer it
correctly rather than confirm it.
\subsection{Scope of This Paper}
\label{sec:scope}
We are explicit about what "framework for quantifying" means here,
since the distinction matters for how the rest of the paper should be
read. The framework --- Computational Entropy, SII, the Mutation
Engine, and PEC --- is a methodology for \emph{measuring} whether
prompt instability predicts inference energy, applicable regardless of
what any single application of it finds. It is not, by itself, a claim
that the underlying cost exists. Section~\ref{sec:results} applies
this methodology at three levels of measurement fidelity and reports
all three results plainly, including that the most rigorous of the
three does not confirm what the least rigorous one suggested. We treat
that outcome as the paper's central empirical contribution rather than
an inconvenience to be minimised, because a framework whose
predictions do not survive contact with real hardware is exactly the
kind of thing the field needs to know about.
\subsection{Contributions}
This paper contributes:
\begin{enumerate}
\item \textbf{Computational Entropy ($S_c$) and the Semantic
Instability Index (SII)}: a theoretically motivated, single-pass,
sub-millisecond proxy (Section~\ref{sec:framework}) for prompt
semantic instability, adapting semantic-entropy
clustering~\cite{kuhn2023semantic} to the input side of inference.
\item \textbf{A nine-class Mutation Engine} (Table~\ref{tab:mutations})
for generating semantically graded, reproducible prompt variants,
including a length-matched control variant used in
Section~\ref{sec:results-stage3}.
\item \textbf{The Prompt Entropy Coefficient (PEC)} and a full
statistical protocol around it (bootstrap confidence intervals, ANOVA
with $\eta^2/\omega^2$, Kruskal--Wallis corroboration, Cohen's $d$,
and a coefficient-sensitivity sweep) for testing SII against measured
energy-per-token.
\item \textbf{A three-tier escalating-fidelity validation}: the same
protocol run under (i) a CPU timing proxy with simulated generation,
(ii) a real, untrained, from-scratch Transformer on real CPU timing,
and (iii) two real, open-weight, instruction-tuned models with real
NVIDIA NVML GPU-energy measurement and an explicit length-matched
control --- to our knowledge the first work to test an LLM
energy-related hypothesis across this specific escalation and report
where it stops replicating.
\item \textbf{An honest negative finding with a diagnosed cause}: the
SII--EPT association present under proxy and untrained-model
measurement is statistically absent under real hardware and real
models, and we identify prompt-length confounding and proxy-formula
linearity as plausible mechanisms by which weaker measurement tiers
manufacture the appearance of an effect.
\end{enumerate}
% =============================================================================
\section{Related Work}
\label{sec:related}
\textbf{Inference energy measurement.} Strubell et
al.~\cite{strubell2019energy} and Patterson et
al.~\cite{patterson2022carbon} established energy/carbon reporting for
training; Samsi et al.~\cite{samsi2023words} extended this to
inference, benchmarking GPU power across batch sizes and model scales
--- our external plausibility check for Stage 1
(Section~\ref{sec:results-stage1}). Faiz et al.'s
\emph{LLMCarbon}~\cite{faiz2024llmcarbon} models end-to-end carbon
footprint and reports baseline per-query energy consistent in order of
magnitude with our own. Bannour et al.~\cite{bannour2021evaluating}
surveyed energy-measurement tooling for NLP, including the
CodeCarbon-style~\cite{lacoste2019quantifying} TDP-proxy approach our
Stage 1 reuses; Desislavov et al.~\cite{desislavov2023trends} found
per-query energy has grown more slowly than parameter count across
model generations, an important caveat when generalising any
single-hardware EPT figure, including ours, across model families.
Wu et al.~\cite{wu2022sustainable} quantified the inference/training
energy crossover, and Chien et al.~\cite{chien2023reducing} and Dodge
et al.~\cite{dodge2022measuring} studied carbon-aware and
region-aware scheduling --- infrastructure-level levers that are
complementary to, not competing with, the prompt-level question we
ask. None of this literature treats prompt content itself as an
independent variable, which is the specific gap this paper addresses.
\textbf{Semantic entropy and prompt engineering.} Our SII is
conceptually downstream of \emph{semantic entropy}: Kuhn et
al.~\cite{kuhn2023semantic} cluster sampled model responses by
bidirectional NLI entailment and compute entropy over the resulting
equivalence classes to detect hallucination, substantially
outperforming token-level entropy; Farquhar et
al.~\cite{farquhar2024detecting} extended this without ground-truth
labels. Malinin and Gales~\cite{malinin2021uncertainty} separately
showed that naive token-level entropy conflates epistemic and
aleatoric uncertainty, motivating cluster-based measures generally.
All of this prior work measures entropy on the \emph{output} side for
quality or hallucination assessment; we adapt the same clustering
construct as $S_c$ (Section~\ref{sec:framework}A) and ask two
different questions: whether a cheap, input-side proxy for it predicts
compute cost (our primary question), and separately
(Section~\ref{sec:results-stage3}) whether that proxy tracks real
output-side $S_c$ directly. Sahoo et al.~\cite{sahoo2024systematic}
catalogued the rapidly growing prompt-engineering literature; none of
the strategies they survey are evaluated for energy cost, underscoring
that this pairing of metrics is, to our knowledge, novel even though
the underlying constructs (semantic entropy, TDP-proxy energy
accounting) are each individually established.
\textbf{Green AI.} Schwartz et al.'s Green AI
agenda~\cite{schwartz2020green} argues efficiency should be a
first-class, reported metric alongside accuracy. This paper is offered
in that spirit, including in reporting a result that complicates
rather than confirms its own premise: we think an honest null,
correctly diagnosed, is more useful to the Green AI literature than a
proxy-measured positive result would have been.
% =============================================================================
\section{Framework: Computational Entropy and SII}
\label{sec:framework}
\subsection{Computational Entropy and Its Motivating Hypothesis}
Given model $\mathcal{M}$ and prompt $x$, partition possible responses
into equivalence clusters $\{C_1,\ldots,C_k\}$ under bidirectional NLI
entailment (two responses share a cluster iff each entails the other),
and let $P(C_i\mid x)$ be the probability mass $\mathcal{M}$ places in
cluster $i$. \emph{Computational Entropy} is
\begin{equation}
  S_c(x) = -\sum_{i=1}^{k} P(C_i \mid x)\,\log_2 P(C_i \mid x).
  \label{eq:sc}
\end{equation}
A coherent, unambiguous prompt concentrates almost all probability
mass on one cluster, giving $S_c\approx0$; a contradictory prompt
spreads mass across clusters that cannot be reconciled, raising $S_c$.
We connect $S_c$ to energy only as a stated hypothesis, not a
derivation. Landauer's Principle~\cite{landauer1961irrev},
experimentally verified for single-bit systems by B\'erut et
al.~\cite{berut2012experimental}, ties irreversible bit-erasure to a
thermodynamic energy floor of $k_BT\ln2$ per bit. Bennett
\cite{bennett1982thermo} showed reversible computation can avoid this
floor in principle, but Transformer decoding~\cite{vaswani2017attention}
--- softmax attention and logit sampling at every step --- is
irreversible in the relevant sense: the pre-commitment distribution
cannot be recovered from the sampled token. \emph{If} the effective
number of such irreversible steps, $N_\text{irrev}(x)$, scales with
$S_c(x)$, then $E\approx N_\text{irrev}(x)\cdot k_BT\ln2$, and higher
semantic instability should raise measured energy. We state plainly
that GPU/CPU dissipation exceeds this Landauer bound by roughly ten
orders of magnitude due to non-equilibrium switching and resistive
losses, so equation (\ref{eq:sc})'s role is to motivate a
\emph{directional} test, not to predict a magnitude, and
Section~\ref{sec:results} is exactly that test, run three times at
increasing rigor.
\subsection{Semantic Instability Index (SII)}
Computing $S_c$ requires multiple generations and pairwise NLI
classification per prompt --- expensive enough to defeat the purpose
of an energy-measurement study if used as the independent variable
directly. We instead define a single-pass proxy:
\begin{equation}
  \mathrm{SII}(x,m,\alpha) = b_m\alpha + \gamma|\Delta F| +
  \delta(1-\mathrm{LD}) + \varepsilon C_x,
  \label{eq:sii}
\end{equation}
where $b_m$ is a stated base score per mutation type $m$
(Table~\ref{tab:mutations}), $\alpha\in[0,1]$ is mutation intensity,
$\Delta F=F(x)-F(x_0)$ is the change in Flesch Reading
Ease~\cite{flesch1948readability} relative to baseline prompt $x_0$,
$\mathrm{LD}$ is lexical diversity (unique-token ratio), and
$C_x=\bar\ell(x)/\max\ell(x)$ is mean-to-max sentence-length ratio.
Each term targets a distinct instability source: $b_m\alpha$ encodes
intended mutation severity, $\Delta F$ readability degradation,
$(1-\mathrm{LD})$ repetitive vocabulary, $C_x$ uneven sentence
structure. $\gamma=0.02,\delta=0.50,\varepsilon=0.30$ were tuned on a
held-out 12-prompt pilot corpus and fixed for all subsequent
measurement, to avoid circularity between calibration and evaluation
data. $b_m$ is a stated design choice --- surface noise scores lowest,
logical contradiction highest (2.50) --- validated only by whether
independently measured energy tracks it, which
Section~\ref{sec:results} tests at three fidelities with different
outcomes, and Section~\ref{sec:results-stage2} subjects to a
coefficient-sensitivity sweep rather than asserting it is correctly
calibrated.
\begin{table}[t]
\caption{Mutation Types, Operations, and SII Base Scores}
\label{tab:mutations}
\centering
\setlength{\tabcolsep}{3.2pt}
\begin{tabular}{@{}llc@{}}
\toprule
\textbf{Type} & \textbf{Operation} & $b_m$ \\
\midrule
\texttt{baseline} & Identity (no change) & 0.00 \\
\texttt{noise\_typo} & Keyboard-adjacency character edits & 0.80 \\
\texttt{noise\_verbose} & Filler-phrase boundary insertion & 1.00 \\
\texttt{formality\_shift} & Register-pair word substitution & 1.20 \\
\texttt{code\_switching} & Parenthetical L2 phrase insertion & 1.40 \\
\texttt{ambiguity\_semantic} & Vague pronoun/quantifier injection & 1.50 \\
\texttt{negation} & Negation insertion & 1.60 \\
\texttt{reordering} & Word/sentence-level shuffle & 1.70 \\
\texttt{ambiguity\_contradiction} & Contradictory clause injection & 2.50 \\
\bottomrule
\end{tabular}
\end{table}
The Mutation Engine applies these nine transformations to each base
prompt with a uniqueness check (re-running with a fallback transform
if a mutation leaves text unchanged --- e.g.\ negation falls back to
inserting a clause-boundary negation phrase when no direct negation
target is found). For the real-model replication
(Section~\ref{sec:results-stage3}) it additionally produces a
length-matched variant per mutation: the mutated text truncated at a
sentence boundary, or padded with neutral filler phrases (never a
literal pad token, which would be out-of-distribution for an
instruction-tuned model), until its token count under the target
model's own tokenizer equals the baseline's exactly.
Algorithm~\ref{alg:sii} gives the SII computation in full.
\begin{algorithm}[t]
\caption{Semantic Instability Index}
\label{alg:sii}
\begin{algorithmic}[1]
\REQUIRE Prompt $x$, baseline $x_0$, mutation type $m$, intensity $\alpha$
\ENSURE $\mathrm{SII}\in[0,5]$
\STATE $F_0\gets\textsc{Flesch}(x_0)$; $F\gets\textsc{Flesch}(x)$; $\Delta F\gets F-F_0$
\STATE $\mathrm{LD}\gets|\{\text{unique tokens in }x\}|/|\text{tokens in }x|$
\STATE $C_x\gets\bar\ell(x)/\max\ell(x)$ \COMMENT{mean/max sentence length}
\STATE $b_m\gets\textsc{SII\_BASE}[m]$ \COMMENT{Table~\ref{tab:mutations}}
\STATE $\mathrm{SII}\gets b_m\alpha+\gamma|\Delta F|+\delta(1-\mathrm{LD})+\varepsilon C_x$
\RETURN $\min(\mathrm{SII},5.0)$
\end{algorithmic}
\end{algorithm}
% =============================================================================
\section{Three-Tier Experimental Protocol}
\label{sec:protocol}
All three tiers share the same 60-prompt corpus (12 each across
factual QA, instruction following, summarisation, code explanation,
and creative writing, so that any effect must survive varying
output-length regimes rather than being an artefact of one task type),
the same nine mutation conditions at $\alpha=0.7$ (chosen after a
pilot comparing $\alpha\in\{0.5,0.7,0.9\}$: $0.5$ under-mutated short
prompts, $0.9$ produced ungrammatical outputs that triggered early-exit
artefacts unrelated to semantic difficulty), and the same statistical
protocol: Spearman's PEC with $10{,}000$-resample bootstrap confidence
intervals, one-way ANOVA with $\eta^2/\omega^2$ and Kruskal--Wallis
corroboration, Cohen's $d$ between baseline and
\texttt{ambiguity\_contradiction}, and (Tier 3 only) a paired Wilcoxon
signed-rank test between raw and length-matched EPT. Sample size
follows an a priori power analysis (Cohen~\cite{cohen1988statistical},
medium effect $f=0.25$, 9 groups, $\alpha=0.05$, $1-\beta=0.95$
indicating $N\approx270$; Tiers 1--2 use $N=540$ for adequate power on
pairwise as well as omnibus comparisons).
The tiers differ only in how text generation and energy are obtained.
\textbf{Tier 1 (proxy)}: energy is $E=P_\text{TDP}f_\text{cpu}
\tau_\text{wall}$, where $f_\text{cpu}=\min(\tau_\text{cpu}/
(\tau_\text{wall}N_\text{cores}),1)$ and $P_\text{TDP}=65$\,W is a
nominal manufacturer-rated constant, not an instantaneous power
reading, plus a synthetic term $E_\text{sim}=\varepsilon_0(n_\text{in}
+n_\text{out})\lambda$ ($\varepsilon_0=0.035$\,J/token) calibrated so
$E+E_\text{sim}$ falls within the range reported
by~\cite{samsi2023words,faiz2024llmcarbon}; generation itself is
simulated via a complexity scalar $\lambda$. \textbf{Tier 2 (untrained,
real-hardware)}: $E_\text{sim}$ is removed entirely; text is generated
by an actual, from-scratch decoder-only Transformer implemented in
NumPy (2 layers, 4 heads, $d_\text{model}=64$, byte-level vocabulary,
KV-cached causal attention verified against an uncached
brute-force reference to $10^{-15}$ precision), with randomly
initialised, untrained weights, executed on this evaluation's real
2-core CPU; wall-clock and process CPU time are measured via
\texttt{time.perf\_counter} and \texttt{resource.getrusage}, not
estimated. Weights are untrained because no pretrained model was
network-reachable in the sandboxed environment used for this tier
(egress to Hugging Face and PyTorch's package index was blocked by a
network allowlist), so generated text is not semantically meaningful
--- this tier isolates the architecture's real compute-timing
behaviour from any semantic-understanding effect. \textbf{Tier 3 (real
models)}: two open-weight, instruction-tuned models, \texttt{phi3.5:3.8b}
(Phi-3.5-mini-instruct, 3.8B parameters) and \texttt{gemma2:2b}
(Gemma-2-2B-it, 2.6B parameters), both Q4\_0-quantised (GGUF) and
served locally via Ollama on a single consumer GPU (NVIDIA RTX~3050
Laptop, 4\,GB VRAM, 30\,W power cap; Intel Core~i5-12500H; Windows~11),
generate real text with natural end-of-sequence stopping (a 96-token
safety ceiling, not a fixed length). Energy is genuine hardware-measured
GPU power: a background thread polls \texttt{nvmlDeviceGetPowerUsage}
throughout each generation call and the resulting trace is integrated
by the trapezoidal rule,
\begin{equation}
  E_\text{real} = \int_{t_0}^{t_1} P(t)\,dt \approx
  \sum_i \tfrac{1}{2}\big(P_i+P_{i+1}\big)\big(t_{i+1}-t_i\big),
  \label{eq:nvml}
\end{equation}
giving joules for the GPU package as a whole (not a per-process
partition --- any concurrent GPU load would be included; CPU energy is
not covered, since Intel RAPL requires Linux root access unavailable
on this Windows machine). Each model's own Hugging Face tokenizer
(\texttt{microsoft/Phi-3.5-mini-instruct}; for Gemma-2, the ungated
mirror \texttt{unsloth/gemma-2-2b-it}, since the canonical repository
requires authenticated licence acceptance) supplies exact token counts
for length matching, with no model weights downloaded from either
repository --- generation runs entirely through the locally served
Ollama GGUF weights. A stratified 15-prompt subset (3 per domain) was
used to keep total run time near an hour per model, giving $N=255$ per
model (135 raw $+$ 120 length-matched; baseline has no length-matched
counterpart by construction).
% =============================================================================
\section{Results}
\label{sec:results}
\subsection{Tier 1: Proxy-Simulated Protocol}
\label{sec:results-stage1}
Table~\ref{tab:ept} shows mean EPT per condition ($N=540$). The rank
order follows the SII scale without inversion, except
\texttt{code\_switching} sits slightly below \texttt{ambiguity\_semantic}
(attributed to compact subword tokenisation of the inserted L2
phrases). PEC $=0.408$ ($p<0.001$, CI $[0.328,0.484]$). One-way ANOVA:
$F(8,531)=12.34$, $p<0.001$, $\eta^2=0.157$ ($\omega^2=0.144$;
Kruskal--Wallis $H=94.2$, $p<0.001$; Levene $p=0.043$, a marginal
homogeneity violation the Kruskal--Wallis corroboration addresses).
\texttt{ambiguity\_contradiction} shows a $+79.0\,\%$ EPT premium over
baseline (Cohen's $d=1.42$, 95\,\% CI $[1.01,1.83]$), not explained by
output token count (Spearman(tokens,EPT)$=0.09$, n.s.) or wall time
alone (partial correlation controlling for wall time, $\rho=0.28$,
$p<0.001$). Baseline EPT (12.4\,mJ/tok) is order-of-magnitude
consistent with~\cite{samsi2023words,faiz2024llmcarbon} --- our only
external check on this proxy's plausibility, not evidence of its
absolute accuracy.
\begin{table}[t]
\caption{Tier 1: Mean EPT per Condition ($n=60$; SD in parentheses)}
\label{tab:ept}
\centering
\setlength{\tabcolsep}{3pt}
\begin{tabular}{@{}lrr@{}}
\toprule
\textbf{Condition} & \textbf{Mean EPT (mJ/tok)} & \textbf{$\Delta\%$ vs.\ baseline} \\
\midrule
\texttt{baseline} & 12.4\,(1.9) & --- \\
\texttt{noise\_typo} & 14.4\,(2.1) & $+16.1$ \\
\texttt{noise\_verbose} & 14.9\,(2.3) & $+20.2$ \\
\texttt{formality\_shift} & 15.6\,(2.4) & $+25.8$ \\
\texttt{code\_switching} & 16.1\,(2.6) & $+29.8$ \\
\texttt{ambiguity\_semantic} & 17.0\,(2.7) & $+37.1$ \\
\texttt{negation} & 17.7\,(2.9) & $+42.7$ \\
\texttt{reordering} & 18.4\,(3.1) & $+48.4$ \\
\texttt{ambiguity\_contradiction} & 22.2\,(3.8) & $+79.0$ \\
\bottomrule
\end{tabular}
\end{table}
\subsection{Tier 2: Untrained Real-Hardware Pilot}
\label{sec:results-stage2}
With $E_\text{sim}$ removed and real CPU timing on an actually
executing (if untrained) Transformer, the direction replicates at
smaller magnitude: PEC $=0.252$ ($p<0.001$, CI $[0.171,0.329]$); ANOVA
$F(8,531)=6.23$, $p<0.001$, $\eta^2=0.086$ ($\omega^2=0.072$;
Kruskal--Wallis $H=85.8$, $p<0.001$; Levene $p=0.653$, no violation).
Cohen's $d$ between baseline and \texttt{ambiguity\_contradiction} is
$0.87$. Critically, input prompt length correlates strongly with EPT
here ($\rho=0.863$, $p<10^{-160}$), and the partial SII--EPT
correlation controlling for input token count collapses to
$-0.037$ --- indistinguishable from zero. Mean input length by
condition confirms a plausible mechanism: \texttt{ambiguity\_contradiction}
and \texttt{noise\_verbose} append text (mean 114.6 and 93.8 tokens
respectively, vs.\ 70.1 at baseline), while \texttt{reordering} only
permutes existing tokens (70.1 tokens, identical to baseline) --- and
\texttt{reordering}'s measured EPT is, correspondingly, the
\emph{lowest} of all nine conditions, an inversion the a priori $b_m$
ranking (Table~\ref{tab:mutations}) does not predict.
A 200-draw coefficient-sensitivity sweep ($\pm30\,\%$ joint
perturbation of $\gamma,\delta,\varepsilon$, and all $b_m$) shows two
different things survive differently. PEC's \emph{direction and
significance} are robust: it stayed positive throughout (mean
$0.238$, SD $0.039$, range $[0.098,0.334]$), and every single-coefficient
perturbation left PEC within $0.005$ of the stated value. The
mutation-by-mutation \emph{ordering}, however, is not: mean rank
agreement between empirical and a priori condition order was only
$\rho=0.29$, occasionally negative ($\rho_\text{min}=-0.18$). In
short, on real but untrained hardware, most of the apparent
SII--EPT association is explainable by mutations changing prompt
length rather than by semantic content once length is held constant
--- a specific, testable confound that Tier 3 controls for directly.
\subsection{Tier 3: Real Instruction-Tuned Models, Real GPU Energy}
\label{sec:results-stage3}
This is the paper's central result. Table~\ref{tab:stage-summary}
compares PEC across all three tiers; Table~\ref{tab:real-pec} gives
the Tier-3 breakdown by model and length-control variant.
\begin{table}[t]
\caption{PEC Across All Three Tiers}
\label{tab:stage-summary}
\centering
\setlength{\tabcolsep}{3pt}
\begin{tabular}{@{}lrrl@{}}
\toprule
\textbf{Tier} & \textbf{PEC} & \textbf{$p$} & \textbf{ANOVA $\eta^2$ ($p$)} \\
\midrule
1: Proxy, simulated gen. & 0.408 & $<0.001$ & 0.157 ($<0.001$) \\
2: Untrained, real timing & 0.252 & $<0.001$ & 0.086 ($<0.001$) \\
3a: Phi-3.5, real NVML & $-0.043$ & 0.499 & 0.019 (0.779) \\
3b: Gemma-2, real NVML & $-0.017$ & 0.784 & 0.036 (0.322) \\
\bottomrule
\end{tabular}
\end{table}
For both Phi-3.5 and Gemma-2, PEC is statistically indistinguishable
from zero, \emph{both before and after length matching}
(Table~\ref{tab:real-pec}; partial correlation controlling for input
tokens likewise near zero: $-0.041$ and $-0.015$ respectively). ANOVA
across the nine conditions is non-significant for both models
(Phi-3.5: $F(8,246)=0.60$, $p=0.78$; Gemma-2: $F(8,246)=1.16$,
$p=0.32$; Kruskal--Wallis corroborates both, $p=0.88$ and $p=0.55$);
Cohen's $d$ between baseline and \texttt{ambiguity\_contradiction} is
negligible for both ($0.083$, $-0.059$) --- for Gemma-2 the
contradiction condition is, if anything, slightly \emph{cheaper} than
baseline. This directly contradicts the Tier-1 ``Contradiction Tax''
finding at this scale and with these models. A 200-draw sensitivity
sweep, identical in method to Tier 2's, confirms the null is stable
rather than fragile: mean PEC stayed at $-0.032$ (Phi-3.5, SD
$0.016$) and $-0.005$ (Gemma-2, SD $0.033$) under the same
$\pm30\,\%$ joint perturbation, with $0\,\%$ of draws for either model
exceeding $\rho=0.30$.
\begin{table}[t]
\caption{Tier 3: SII--EPT Correlation, Raw vs.\ Length-Matched}
\label{tab:real-pec}
\centering
\setlength{\tabcolsep}{3pt}
\begin{tabular}{@{}lrr@{}}
\toprule
\textbf{Model / subset} & \textbf{PEC ($\rho$)} & \textbf{$p$} \\
\midrule
Phi-3.5, raw ($n=135$) & $-0.050$ & 0.562 \\
Phi-3.5, length-matched ($n=120$) & $-0.045$ & 0.602 \\
Gemma-2, raw ($n=135$) & $-0.066$ & 0.445 \\
Gemma-2, length-matched ($n=120$) & $+0.030$ & 0.734 \\
\bottomrule
\end{tabular}
\end{table}
The length-matched control was informative on its own terms even
though PEC stayed null. A paired Wilcoxon signed-rank test between raw
and length-matched EPT for identical (prompt, mutation) pairs is
significant for Phi-3.5 ($n=120$, $W=2561$, $p=0.005$; raw EPT higher
by a mean $29.9$\,mJ/token) but not for Gemma-2 ($W=3240$, $p=0.307$;
mean difference $7.2$\,mJ/token). Prompt length can measurably move
EPT for at least one model --- just not in a way that produces any
corresponding aggregate SII--EPT correlation for either model, since
that correlation was already null before length was controlled.
We additionally ran the $S_c$-validation pilot ($k=5$ samples/prompt,
temperature 0.7, bidirectional-NLI clustering following Kuhn et
al.~\cite{kuhn2023semantic}, \texttt{cross-encoder/nli-deberta-v3-xsmall},
threshold 0.5, on a 10-prompt $\times$ 6-condition subset, $n=60$ per
model) that neither earlier tier could execute meaningfully (Tier 2's
untrained model produces incoherent text that cannot be meaningfully
clustered). Results disagree by model: Spearman(SII,$S_c)=+0.345$
($p=0.007$, CI $[0.094,0.553]$) for Gemma-2, in the theoretically
predicted direction, but $-0.173$ ($p=0.187$, CI $[-0.419,0.084]$,
n.s.) for Phi-3.5. We flag a specific measurement concern rather than
taking either number at face value: $50\,\%$ of Phi-3.5's records and
$42\,\%$ of Gemma-2's hit the maximum possible cluster count ($k$
distinct clusters from $k$ samples --- the entailment check never
judged any two responses equivalent, even near-paraphrase answers to
plain factual questions). This ceiling suggests the single-representative,
threshold-0.5 clustering rule may over-fragment responses broadly,
inflating $S_c$ toward near-uniform high values regardless of
condition, which could itself explain the cross-model disagreement
rather than a genuine architectural difference. We report both
numbers but treat neither as settling whether SII tracks real semantic
entropy; Section~\ref{sec:limits} lists the specific fix needed before
this pilot's numbers should be trusted on their own.
% =============================================================================
\section{Discussion}
\label{sec:discussion}
\subsection{Why a Proxy Would Show an Effect That Real Hardware Does Not}
Two mechanisms plausibly explain the gap between Tiers 1--2 and Tier
3. First, length confounding: Tier 2 showed directly that an untrained
model's real compute cost tracks context length, and several mutations
(verbose, contradiction) happen to lengthen prompts; Tier 1's synthetic
term is also linear in token count, so any length--condition
association mechanically produces a PEC signal that need not involve
semantics at all. Tier 3's length-matched control was designed to rule
this out for real models, and did: PEC stayed null even after
equalising length, meaning the Tier-1/2 signal is not simply ``the
same real effect, just smaller.'' Second, representational capacity:
a trained model's decoding step performs the same fixed-size matrix
operations (KV-cached attention over the current context, a
feed-forward block) regardless of whether the input content is
coherent or contradictory --- there is no computational hook for
``the model is confused'' in standard causal decoding, whereas an
untrained architecture's timing noise floor is more exposed to
context-length effects precisely because no learned representation is
smoothing it out. We did not instrument attention activations directly
in any tier, so this account is offered as a plausible explanation for
the pattern in Table~\ref{tab:stage-summary}, not a confirmed
mechanism, and is a natural target for the follow-on work in
Section~\ref{sec:conclusion}.
\subsection{Threats to Validity}
\label{sec:limits}
\emph{Construct validity of energy measurement.} Tiers 1--2 use a
nominal TDP or CPU-timing proxy, not an instantaneous power reading;
Tier 3 uses real NVML power but integrates whole-GPU-package draw (not
a per-process partition, so any concurrent GPU load would inflate a
call's measured energy) and does not cover CPU energy at all, since
Intel RAPL requires Linux root access unavailable on any machine used
in this study. \emph{Statistical power at Tier 3.} $N=255$ per model
(15 stratified prompts, vs.\ 60/$N=540$ at Tiers 1--2) keeps total run
time near an hour; Gemma-2's within-condition standard deviations are
comparable in magnitude to its condition means, so this is a
well-supported null at this scale rather than an exhaustively powered
one --- though the point estimates already sit at zero rather than
near a threshold additional power would likely push past.
\emph{Generation capping.} $90\,\%$ (Phi-3.5) and $79\,\%$ (Gemma-2)
of Tier-3 completions hit the 96-token safety ceiling rather than
stopping at end-of-sequence, limiting how much natural output-length
variation this tier could observe. \emph{Quantisation.} Both Tier-3
models are Q4\_0-quantised GGUF weights served via Ollama, not
full-precision references; any null or positive finding describes this
deployment configuration specifically. \emph{Two models only.}
Phi-3.5 and Gemma-2 disagreeing on the $S_c$ pilot
(Section~\ref{sec:results-stage3}) is itself evidence that
model-to-model variance can be large, and the clustering-ceiling issue
noted there means that disagreement should not be over-interpreted
before the clustering rule is revised --- comparing each candidate
response against all existing cluster members (majority-entailment)
rather than only the first, and calibrating the entailment threshold
against a corpus of known-equivalent and known-distinct response
pairs, are the two concrete fixes we would prioritise. \emph{External
validity.} All three tiers use an author-constructed 60-prompt corpus
across five general domains; results have not been tested on natural,
user-authored prompt distributions (e.g.\ ShareGPT, Alpaca), and
mutation base scores $b_m$ were calibrated on English text only.
\subsection{What This Does and Does Not License}
Pending broader replication, we consider the framework itself ---
SII as a sub-millisecond, single-pass score; the Mutation Engine as a
reproducible perturbation protocol, including its length-matched
control; PEC as a rank-correlation statistic with bootstrap CIs and a
coefficient-sensitivity sweep built in --- a reusable methodological
contribution independent of the sign of any particular application's
result. What the evidence in this paper does not license is a claim
that prompt semantic instability is a demonstrated, exploitable energy
cost in real deployed models: contradiction screening, verbosity
trimming, or any other ``Green Prompt Engineering'' intervention
motivated by an SII--EPT relationship should be treated as unproven
until shown on real hardware, for the specific model and quantisation
in question, with a length-matched control --- exactly the standard
Tier 3 applied here and did not meet. More broadly, we think this
result is a useful caution for the growing body of LLM-energy research
that relies on CPU-timing or TDP-style proxies without real-hardware
validation: Tiers 1--2 of this very paper would, read alone, have
supported a confident positive claim that Tier 3 does not sustain.
% =============================================================================
\section{Conclusion and Future Work}
\label{sec:conclusion}
We proposed a framework --- Computational Entropy ($S_c$), its
single-pass proxy SII, the Mutation Engine, and the Prompt Entropy
Coefficient (PEC) --- for testing whether prompt semantic instability
predicts LLM inference energy, and evaluated it at three levels of
measurement fidelity rather than one. A CPU timing-proxy protocol with
simulated generation shows a significant association (PEC$=0.408$); a
real, untrained Transformer measured with real CPU timing reproduces
the direction at smaller magnitude but attributes most of it to prompt
length rather than semantics (partial correlation $\approx0$ once
length is controlled); two real, instruction-tuned models measured
with real NVIDIA NVML GPU energy and an explicit length-matched
control show no SII--EPT association at all, robust to a
coefficient-sensitivity sweep. We treat this last result as the
paper's main empirical finding: on the strongest evidence gathered
here, prompt semantic instability is not a demonstrated correlate of
per-token inference energy in real, deployed models, even though
progressively weaker measurement tiers suggested otherwise. The
practical implication is methodological as much as substantive: energy
claims about prompt properties, in this work or elsewhere, should be
treated with caution until shown on real hardware with real models and
an explicit length-matched control, because proxy and untrained-model
measurements can manufacture an association that does not survive that
test. Two narrower findings keep the underlying question open rather
than closed --- Gemma-2's significant SII--$S_c$ correlation and
Phi-3.5's significant raw-versus-length-matched energy difference ---
and motivate rather than settle further work.
Priority follow-on work, in order: (1) RAPL- or NVML-based energy
accounting on a broader sweep of model families and parameter scales
than the two tested here, given how much they already disagreed; (2)
a revised $S_c$-clustering protocol addressing the ceiling effect
identified in Section~\ref{sec:results-stage3}, before re-running the
semantic-entropy validation pilot; (3) a larger Tier-3 sample to
tighten confidence intervals around what are currently statistically
null, not precisely small-positive, results; and (4) extension to
multi-turn conversation settings, where context entropy may accumulate
across turns and interact with KV-cache reuse in ways none of the
three tiers here capture.
\subsection*{Availability}
The 60-prompt corpus, Mutation Engine, SII implementation, all three
tiers' measurement scripts, the NumPy Transformer used in Tier 2, the
NVML energy-sampling code and Ollama backend used in Tier 3, and the
complete raw measurement logs underlying every number reported in this
paper are retained by the authors and available for independent
verification on request, to allow full reproduction of
Sections~\ref{sec:results}--\ref{sec:discussion} without re-deriving
any intermediate result from scratch.
% =============================================================================
\balance
\begin{thebibliography}{00}
\bibitem{strubell2019energy}
E.~Strubell, A.~Ganesh, and A.~McCallum,
``Energy and policy considerations for deep learning in NLP,''
in \textit{Proc.\ 57th Annu.\ Meet.\ Assoc.\ Comput.\ Linguist.\ (ACL)},
Florence, Italy, Jul.~2019, pp.~3645--3650.
DOI:~10.18653/v1/P19-1355.
\bibitem{patterson2022carbon}
D.~Patterson et al.,
``Carbon emissions and large neural network training,''
\textit{Commun.\ ACM}, vol.~65, no.~6, pp.~35--35, Jun.~2022.
DOI:~10.1145/3522589.
\bibitem{wu2022sustainable}
C.-J.~Wu et al.,
``Sustainable AI: Environmental implications, challenges,
and opportunities,''
in \textit{Proc.\ 5th Conf.\ Mach.\ Learn.\ Syst.\ (MLSys)},
Santa Clara, CA, USA, 2022.
arXiv:2111.00364.
\bibitem{samsi2023words}
S.~Samsi et al.,
``From words to watts: Benchmarking the energy costs of large
language model inference,''
in \textit{Proc.\ IEEE High Perform.\ Extreme Comput.\ Conf.\
(HPEC)}, Dedham, MA, USA, Sep.~2023, pp.~1--9.
DOI:~10.1109/HPEC58863.2023.10363447.
\bibitem{bannour2021evaluating}
N.~Bannour, S.~Ghannay, A.~N\'ev\'eol, and A.-L.~Ligozat,
``Evaluating the carbon footprint of NLP methods: A survey
and analysis of existing tools,''
in \textit{Proc.\ 2nd Workshop SustaiNLP (EMNLP)},
Nov.~2021, pp.~11--21.
DOI:~10.18653/v1/2021.sustainlp-1.2.
\bibitem{lacoste2019quantifying}
A.~Lacoste, A.~Luccioni, V.~Schmidt, and T.~Dandres,
``Quantifying the carbon emissions of machine learning,''
in \textit{Workshop Tackling Climate Change with AI (NeurIPS)},
Vancouver, Canada, Dec.~2019.
arXiv:1910.09700.
\bibitem{desislavov2023trends}
R.~Desislavov, F.~Mart\'inez-Plumed, and J.~Hern\'andez-Orallo,
``Trends in AI inference energy consumption: Beyond the
performance-vs-parameter laws of deep learning,''
\textit{Sustain.\ Comput.: Inform.\ Syst.}, vol.~38,
art.~100857, Apr.~2023.
DOI:~10.1016/j.suscom.2023.100857.
\bibitem{faiz2024llmcarbon}
A.~Faiz, S.~Kaneda, R.~Wang, R.~Osi, P.~Sharma, F.~Chen, and
L.~Jiang,
``LLMCarbon: Modeling the end-to-end carbon footprint of large
language models,''
in \textit{Proc.\ 12th Int.\ Conf.\ Learn.\ Represent.\ (ICLR)},
Vienna, Austria, May~2024.
arXiv:2309.14393.
\bibitem{chien2023reducing}
A.~A.~Chien, L.~Lin, H.~Nguyen, V.~Rao, T.~Sharma, and
R.~Wijayawardana,
``Reducing the carbon impact of generative AI inference (today
and in 2035),''
in \textit{Proc.\ 2nd Workshop Sustain.\ Comput.\ Syst.\
(HotCarbon)}, Boston, MA, USA, Jul.~2023.
DOI:~10.1145/3604930.3605705.
\bibitem{dodge2022measuring}
J.~Dodge et al.,
``Measuring the carbon intensity of AI in cloud instances,''
in \textit{Proc.\ 2022 ACM Conf.\ Fairness, Accountability,
Transparency (FAccT)}, Seoul, Korea, Jun.~2022,
pp.~1877--1894.
DOI:~10.1145/3531146.3533234.
\bibitem{sahoo2024systematic}
P.~Sahoo, A.~K.~Singh, S.~Saha, V.~Jain, S.~Mondal, and
A.~Chadha,
``A systematic survey of prompt engineering in large language
models: Techniques and applications,''
\textit{arXiv preprint}, 2024.
arXiv:2402.07927.
\bibitem{landauer1961irrev}
R.~Landauer,
``Irreversibility and heat generation in the computing process,''
\textit{IBM J.~Res.\ Dev.}, vol.~5, no.~3,
pp.~183--191, Jul.~1961.
DOI:~10.1147/rd.53.0183.
\bibitem{berut2012experimental}
A.~B\'erut, A.~Arakelyan, A.~Petrosyan, S.~Ciliberto,
R.~Dillenschneider, and E.~Lutz,
``Experimental verification of Landauer's principle linking
information and thermodynamics,''
\textit{Nature}, vol.~483, no.~7388, pp.~187--189, Mar.~2012.
DOI:~10.1038/nature10872.
\bibitem{bennett1982thermo}
C.~H.~Bennett,
``The thermodynamics of computation---a review,''
\textit{Int.~J.~Theor.\ Phys.}, vol.~21, nos.~12,
pp.~905--940, 1982.
DOI:~10.1007/BF02084158.
\bibitem{vaswani2017attention}
A.~Vaswani et al.,
``Attention is all you need,''
in \textit{Adv.\ Neural Inf.\ Process.\ Syst.\ (NeurIPS)},
vol.~30, Long Beach, CA, USA, Dec.~2017, pp.~5998--6008.
arXiv:1706.03762.
\bibitem{kuhn2023semantic}
L.~Kuhn, Y.~Gal, and S.~Farquhar,
``Semantic uncertainty: Linguistic invariances for uncertainty
estimation in natural language generation,''
in \textit{Proc.\ 11th Int.\ Conf.\ Learn.\ Represent.\ (ICLR)},
Kigali, Rwanda, May~2023.
arXiv:2302.09664.
\bibitem{farquhar2024detecting}
S.~Farquhar, J.~Kossen, L.~Kuhn, and Y.~Gal,
``Detecting hallucinations in large language models using
semantic entropy,''
\textit{Nature}, vol.~630, no.~8017, pp.~625--630, Jun.~2024.
DOI:~10.1038/s41586-024-07421-0.
\bibitem{malinin2021uncertainty}
A.~Malinin and M.~Gales,
``Uncertainty estimation in autoregressive structured
prediction,''
in \textit{Proc.\ 9th Int.\ Conf.\ Learn.\ Represent.\ (ICLR)},
May~2021.
arXiv:2002.07650.
\bibitem{schwartz2020green}
R.~Schwartz, J.~Dodge, N.~A.~Smith, and O.~Etzioni,
``Green AI,''
\textit{Commun.\ ACM}, vol.~63, no.~12, pp.~54--63, Dec.~2020.
DOI:~10.1145/3381831.
\bibitem{flesch1948readability}
R.~Flesch,
``A new readability yardstick,''
\textit{J.~Appl.\ Psychol.}, vol.~32, no.~3, pp.~221--233,
1948.
DOI:~10.1037/h0057532.
\bibitem{cohen1988statistical}
J.~Cohen,
\textit{Statistical Power Analysis for the Behavioral Sciences},
2nd~ed. Hillsdale, NJ, USA: Lawrence Erlbaum Associates, 1988.
\end{thebibliography}
\end{document}
